\chapter{Tangential Jacobians and weighted coarea}\label{sec:weighted-coarea}
Integration along fibers is used both to cut currents and to select level sets of controlled perimeter. The relevant derivative on a rectifiable set is tangential; ambient differentiability at almost every ambient point does not by itself define a derivative on that set. We follow \cite{Sim18}, Chapter 3, Section 2, pp.~92--94.

\begin{proposition}[Weighted Euclidean area formula]\label{prop:weighted-area}
Under the hypotheses of Theorem \ref{thm:area-formula}, let $h:V\to[0,\infty]$ be measurable. Then
\[
 \int_V h(x)J_mDf(x)\,dx
 =\int_{\mathbb R^n}\sum_{x\in f^{-1}(y)}h(x)\,d\mathcal H^m(y).
\]
The sum means integration against counting measure, and both sides may be infinite. If $h$ is signed, the same identity holds when the left side with $|h|$ is finite.
\end{proposition}
\begin{proof}
For $h=\chi_A$ this is the area formula. Additivity gives the identity for nonnegative simple functions; monotone convergence gives the general nonnegative case. Apply it to the positive and negative parts for the signed statement. This is the extension in \cite{Sim18}, Chapter 3, equation (2.5), specialized to a Euclidean domain. Integrability with the Jacobian, rather than unweighted integrability alone, is the relevant condition.
\end{proof}

\begin{definition}[Tangential differential and coarea Jacobian]\label{def:tangential-coarea-jacobian}
Let $M\subset U\subset\mathbb R^N$ be countably $m$-rectifiable, $\mathcal H^m$-measurable, and of locally finite $\mathcal H^m$ measure. Let $f:U\to\mathbb R^k$ be locally Lipschitz, where $1\leq k<m$. On countably many $C^1$ manifold pieces covering $M$ up to a null set, differentiate the restrictions of $f$ in manifold charts. Their differentials agree on overlaps almost everywhere, defining
\[
 d^Mf_x:T_xM\to\mathbb R^k,\qquad
 J^M_kf(x)=\sqrt{\det\bigl(d^Mf_x(d^Mf_x)^*\bigr)}.
\]
Set the Jacobian to zero on the exceptional null set. If $k=1$, the tangential gradient represents the differential and $J^M_1f=|\nabla^Mf|$. The existence of the manifold decomposition and the almost-everywhere independence are rectifiability results recorded in \cite{Sim18}, Chapter 3, Definition 2.1 and equations (2.2), (2.7); their local proof chain is given by Theorems \ref{thm:rectifiable-c1-decomposition} and \ref{thm:rectifiable-tangential-differential}, with Whitney extension and the Euclidean area formula retained as inputs. At a point where the ambient differential exists, the tangential differential is its restriction to $T_xM$.
\end{definition}

\begin{theorem}[Coarea on a rectifiable set]\label{thm:rectifiable-coarea}
With $M,f,m,k$ as in Definition \ref{def:tangential-coarea-jacobian}, for every $\mathcal H^m$-measurable $A\subset M$,
\[
 \int_A J^M_k f\,d\mathcal H^m
 =\int_{\mathbb R^k}\mathcal H^{m-k}(A\cap f^{-1}(y))\,dy.
\]
The fiber measure is measurable and the equality allows infinite values. For almost every $y$, the fiber $M\cap f^{-1}(y)$ is countably $(m-k)$-rectifiable and has locally finite $\mathcal H^{m-k}$ measure in $U$.

The coarea and rectifiability assertions are \cite{Sim18}, Chapter 3, equation (2.8) and Remark 2.10(2), pp.~93--94. Their local chart proof, including the fiber tangent geometry, is now supplied by Theorem \ref{thm:rectifiable-coarea-chart-proof}. The Euclidean area/coarea formulas now have local proofs; Whitney extension remains an explicit external input.
\end{theorem}
\begin{proof}
Use Theorem \ref{thm:rectifiable-coarea-chart-proof} with the nonnegative weight $h=\chi_A$. Its integral identity is exactly the displayed formula, and its representative argument supplies the measurable fiber measure. The same theorem gives rectifiability and local finite measure of almost every fiber of $M$. The proof there constructs the geometry from $C^1$ approximation and regular coordinate fibers, uses the chart Jacobian identity to pass from Euclidean coarea to the surface, and handles the omitted Hausdorff-null set through the separately proved null-fiber lemma. Thus no fiber smoothness or ambient differentiability on all of $M$ is being assumed.
\end{proof}

\begin{corollary}[Weighted rectifiable coarea]\label{cor:weighted-rectifiable-coarea}
For nonnegative measurable $g$ on $M$,
\[
 \int_M gJ^M_kf\,d\mathcal H^m
 =\int_{\mathbb R^k}\left(\int_{M\cap f^{-1}(y)}g\,d\mathcal H^{m-k}\right)dy.
\]
The identity also holds for signed $g$ with $\int_M|g|J^M_kf<\infty$. In particular, if $C=\{x:J^M_kf(x)=0\}$, then $\mathcal H^{m-k}(C\cap f^{-1}(y))=0$ for almost every $y$; no assertion that $f(C)$ has measure zero is made.
\end{corollary}
\begin{proof}
Approximate a nonnegative function by increasing simple functions and apply Theorem \ref{thm:rectifiable-coarea} and monotone convergence on both sides. Positive and negative parts give the signed statement. Set $g=\chi_C$ for the last assertion. This is \cite{Sim18}, Chapter 3, equation (2.9).
\end{proof}

\begin{proposition}[Selecting a radial slice of controlled area]\label{prop:rectifiable-radial-slice}
Let $m\geq2$, let $M$ satisfy the rectifiability and local-finiteness hypotheses above, and suppose $\overline{B_b(a)}\subset U$, $0<a_0<b$. Then
\[
 \int_{a_0}^b\mathcal H^{m-1}(M\cap\partial B_r(a))\,dr
 \leq\mathcal H^m(M\cap(B_b(a)\setminus\overline{B_{a_0}(a)})).
\]
Consequently, for every $\varepsilon>0$ there is a radius $r\in(a_0,b)$, outside any prescribed Lebesgue-null exceptional set, such that the slice measure is at most the right side divided by $b-a_0$, plus $\varepsilon$.
\end{proposition}
\begin{proof}
Apply weighted coarea to the distance $\rho(x)=|x-a|$ with the annulus indicator. Its tangential gradient has norm at most one. The average bound implies the existence assertion; otherwise integrating the strict opposite inequality outside the exceptional set would contradict the bound. The exceptional-set clause lets this selection be combined with the almost-everywhere conclusions of the current slicing theorem, Theorem \ref{thm:integral-current-slicing}.
\end{proof}

\begin{example}[Projection recovers Fubini]\label{ex:coarea-fubini}
On $M=\mathbb R^m$, let $f(x_1,\ldots,x_m)=(x_1,\ldots,x_k)$. Then $d^Mf(d^Mf)^*=I_k$ and $J^M_kf=1$. The fibers are affine copies of $\mathbb R^{m-k}$, so weighted coarea becomes Fubini--Tonelli. This is \cite{Sim18}, Chapter 2, Remark 7.4(2).
\end{example}

\begin{example}[An ambient exceptional set can contain the whole surface]\label{ex:tangential-not-ambient}
For $M=\mathbb R^m\times\{0\}\subset\mathbb R^{m+1}$ and $f(x,z)=|z|$, the ambient function is Lipschitz but is not differentiable at any point of $M$. Its restriction to $M$ is constant and its tangential gradient is zero everywhere. Thus assigning an arbitrary ambient derivative on the ambient Lebesgue-null exceptional set is not a definition of the tangential derivative. Differentiation in manifold charts is essential.
\end{example}
